
\documentclass[10pt]{article}
\usepackage[utf8]{inputenc}
\usepackage[T1]{fontenc}
\usepackage{amsmath, amssymb, xcolor, geometry, enumitem, multicol, tcolorbox}

\definecolor{mainblue}{RGB}{0, 102, 204}
\geometry{a4paper, margin=1.2cm, top=1cm, bottom=3cm, footskip=1.2cm}

\begin{document}
\enlargethispage{2.5cm}

% Modern Header
\begin{tcolorbox}[colback=mainblue!5, colframe=mainblue, sharp corners, boxrule=0.5pt]
    \small \textbf{Unit:} undefined \hfill \textbf{Name:} \rule{3cm}{0.4pt} \\
    \small \textbf{Topic:} Variables and Expressions / Order of Operations \hfill \textbf{Date:} \rule{2cm}{0.4pt} \\
    \centering \textbf{\large Classwork}
\end{tcolorbox}

\vspace{0.2cm}

% MCQ Section
\begin{multicols}{2}
\begin{enumerate}[label=\textbf{Q\arabic*.}, leftmargin=*, itemsep=6pt, parsep=0pt]

  \item \small What is the variable in the expression $2x + 3$?
  \begin{enumerate}[label=(\Alph*), leftmargin=0.6cm, nosep, itemsep=2pt]
    \item 2
    \item 3
    \item x
    \item 2x
    \item 3x
  \end{enumerate}

  \item \small Identify the coefficient in the term $4y$
  \begin{enumerate}[label=(\Alph*), leftmargin=0.6cm, nosep, itemsep=2pt]
    \item y
    \item 4
    \item 4y
    \item 1
    \item 0
  \end{enumerate}

  \item \small What operation is represented by the symbol $+$ in $2 + 5$?
  \begin{enumerate}[label=(\Alph*), leftmargin=0.6cm, nosep, itemsep=2pt]
    \item Subtraction
    \item Addition
    \item Multiplication
    \item Division
    \item Exponentiation
  \end{enumerate}

  \item \small In the expression $3 cdot 2$, what is the operation being performed?
  \begin{enumerate}[label=(\Alph*), leftmargin=0.6cm, nosep, itemsep=2pt]
    \item Addition
    \item Subtraction
    \item Multiplication
    \item Division
    \item Exponentiation
  \end{enumerate}

  \item \small What is the constant term in $2x - 4$?
  \begin{enumerate}[label=(\Alph*), leftmargin=0.6cm, nosep, itemsep=2pt]
    \item 2x
    \item 2
    \item -4
    \item x
    \item 0
  \end{enumerate}

  \item \small Identify the operator in the expression $7 - 2$
  \begin{enumerate}[label=(\Alph*), leftmargin=0.6cm, nosep, itemsep=2pt]
    \item +
    \item -
    \item $cdot$
    \item $/$
    \item $^$
  \end{enumerate}

  \item \small What is the term $5x$ an example of?
  \begin{enumerate}[label=(\Alph*), leftmargin=0.6cm, nosep, itemsep=2pt]
    \item Constant
    \item Variable
    \item Coefficient
    \item Monomial
    \item Polynomial
  \end{enumerate}

  \item \small In the expression $9 + 1$, what operation is being performed?
  \begin{enumerate}[label=(\Alph*), leftmargin=0.6cm, nosep, itemsep=2pt]
    \item Subtraction
    \item Addition
    \item Multiplication
    \item Division
    \item Exponentiation
  \end{enumerate}

\end{enumerate}
\end{multicols}

\vspace{-0.2cm}
\hrule
\vspace{0.2cm}
\textbf{\small Open Ended Questions (FRQ)}
\begin{enumerate}[label=\textbf{Q\arabic*.}, start=9, itemsep=5pt, topsep=0pt]

  \item \small Explain, using examples, the difference between a variable and a constant in an algebraic expression. Provide a simple expression that includes both. \vspace{2cm}

  \item \small Draw a simple diagram to illustrate the concept of the order of operations. Use the expression $3 + 2 cdot 4$ as an example and explain how it should be evaluated step by step. \vspace{2cm}

\end{enumerate}

\vfill

% Chef's Tips Footer
\begin{tcolorbox}[colback=blue!5, colframe=mainblue!50, title=\small \textbf{Chef's Tips}, fonttitle=\bfseries, bottom=1mm]
\begin{itemize}[noitemsep, topsep=2pt, leftmargin=0.5cm]
  \item \scriptsize Ensure students understand the basic definitions of variables, constants, and coefficients before proceeding to more complex topics.\item \scriptsize Use visual aids and simple examples to explain the order of operations to help students grasp the concept effectively.\item \scriptsize Provide ample practice opportunities for students to identify and apply the order of operations in various expressions.
\end{itemize}
\end{tcolorbox}

\end{document}
  